\documentclass[a4paper,12pt]{article}

\usepackage[T2A]{fontenc}
\usepackage[utf8]{inputenc}
\usepackage[russian]{babel}

\usepackage{amsmath, amssymb, amsthm, mathtools}
\usepackage{helvet}
\renewcommand{\familydefault}{\sfdefault}
\usepackage{icomma}
\usepackage{geometry}
\usepackage{microtype}
\usepackage{tikz}
\usepackage{tkz-euclide}
\usepackage{pgfplots}
\pgfplotsset{compat=1.18}
\usepackage{tabularx, booktabs, longtable}
\usepackage{enumitem}
\usepackage{hyperref}
\usepackage{parskip}
\usepackage{fancyhdr}
\usepackage{setspace}
\usepackage{xcolor}

\geometry{left=18mm,right=18mm,top=18mm,bottom=20mm}
\onehalfspacing
\sloppy

\definecolor{accent}{HTML}{008080}
\definecolor{darkaccent}{HTML}{004D4D}
\definecolor{textgray}{HTML}{333333}
\definecolor{softgray}{HTML}{F4F6F6}

\hypersetup{
    colorlinks=true,
    linkcolor=darkaccent,
    urlcolor=darkaccent
}

\pagestyle{fancy}
\fancyhf{}
\lhead{\textcolor{textgray}{Чек-лист: задачи на движение}}
\rhead{\textcolor{textgray}{ОГЭ}}
\cfoot{\thepage}
\renewcommand{\headrulewidth}{0.4pt}
\renewcommand{\headrule}{\hbox to\headwidth{\color{accent}\leaders\hrule height \headrulewidth\hfill}}

\setlist[itemize]{leftmargin=*, itemsep=2pt, topsep=3pt}
\setlist[enumerate]{leftmargin=*, itemsep=4pt, topsep=4pt}

\newcommand{\important}[1]{\textcolor{darkaccent}{\textbf{#1}}}
\newcommand{\kmh}{\text{км/ч}}
\newcommand{\km}{\text{км}}
\newcommand{\ch}{\text{ч}}
\newcommand{\m}{\text{м}}

\newcommand{\motionlineAB}[4]{
\begin{center}
\begin{tikzpicture}[scale=0.95]
    \draw[line width=0.8pt] (0,0) -- (9,0);
    \fill[accent] (0,0) circle (2pt);
    \fill[accent] (9,0) circle (2pt);
    \node[below] at (0,0) {$A$};
    \node[below] at (9,0) {$B$};
    \node[above] at (4.5,0.1) {#1};
    \draw[->, line width=1pt, color=darkaccent] (0.6,0.45) -- (3.5,0.45);
    \draw[->, line width=1pt, color=darkaccent] (8.4,0.45) -- (5.5,0.45);
    \node[above] at (2.05,0.45) {#2};
    \node[above] at (6.95,0.45) {#3};
    \node[below] at (4.5,-0.35) {#4};
\end{tikzpicture}
\end{center}
}

\newcommand{\catchupline}[4]{
\begin{center}
\begin{tikzpicture}[scale=0.95]
    \draw[line width=0.8pt] (0,0) -- (9,0);
    \fill[accent] (0,0) circle (2pt);
    \fill[accent] (3,0) circle (2pt);
    \fill[accent] (9,0) circle (2pt);
    \node[below] at (0,0) {$A$};
    \node[below] at (3,0) {$B$};
    \node[below] at (9,0) {$C$};
    \node[above] at (1.5,0.1) {#1};
    \draw[->, line width=1pt, color=darkaccent] (0.5,0.45) -- (4.2,0.45);
    \draw[->, line width=1pt, color=darkaccent] (3.5,0.85) -- (6.8,0.85);
    \node[above] at (2.35,0.45) {#2};
    \node[above] at (5.15,0.85) {#3};
    \node[below] at (6.2,-0.35) {#4};
\end{tikzpicture}
\end{center}
}

\begin{document}

\begin{titlepage}
\thispagestyle{empty}
\begin{center}
\vspace*{25mm}

{\Large \textcolor{accent}{\textbf{Чек-лист}}}

\vspace{8mm}

{\Huge \textbf{Задачи на движение}}

\vspace{4mm}

{\Large движение навстречу, вдогонку и движение туда-обратно}

\vspace{22mm}

{\large Лёвин Артём Александрович\\[1mm]
эксклюзивно для Софии Кальней}

\vspace{12mm}

{\large \today}

\vfill

\begin{tikzpicture}
    \draw[color=accent, line width=1.2pt]
    (0,0) .. controls (2,0.8) and (4,-0.8) .. (6,0)
          .. controls (8,0.8) and (10,-0.8) .. (12,0);
    \draw[color=darkaccent, line width=0.5pt]
    (1,-0.35) .. controls (3,0.25) and (5,-0.95) .. (7,-0.35)
          .. controls (9,0.25) and (10.5,-0.75) .. (11.5,-0.35);
\end{tikzpicture}

\vspace{8mm}
\end{center}
\end{titlepage}

\tableofcontents
\newpage

\section{Главная идея задач на движение}

\important{Базовая формула:}
\[
S = vt.
\]

Здесь:
\[
S \text{ --- расстояние}, \qquad
v \text{ --- скорость}, \qquad
t \text{ --- время}.
\]

Из этой формулы получаем ещё две:
\[
v=\frac{S}{t}, \qquad t=\frac{S}{v}.
\]

\important{Главное правило:} перед составлением уравнения нужно понять, какие величины складываются, какие равны, а какие отличаются на заданную величину.

\section{Единицы измерения}

В задачах ОГЭ часто смешиваются километры, метры, часы и минуты.

\begin{center}
\renewcommand{\arraystretch}{1.35}
\begin{tabularx}{0.85\linewidth}{X X}
\toprule
\textbf{Величина} & \textbf{Как переводить} \\
\midrule
$300\ \m$ & $0{,}3\ \km$ \\
$600\ \m$ & $0{,}6\ \km$ \\
$15$ минут & $\dfrac{15}{60}=0{,}25\ \ch$ \\
$30$ минут & $\dfrac{30}{60}=0{,}5\ \ch$ \\
$45$ минут & $\dfrac{45}{60}=0{,}75\ \ch$ \\
\bottomrule
\end{tabularx}
\end{center}

\important{Типичная ошибка:} нельзя складывать километры с метрами или часы с минутами без перевода к одним единицам.

\section{Движение навстречу друг другу}

\subsection{Ситуация}

Два объекта движутся из разных пунктов навстречу друг другу.

\motionlineAB{$S$}{$v_1$}{$v_2$}{движение навстречу}

Если они выехали одновременно и встретились через $t$ часов, то первый прошёл $v_1t$, второй прошёл $v_2t$, а вместе они прошли всё расстояние:
\[
v_1t+v_2t=S.
\]

Можно записать короче:
\[
(v_1+v_2)t=S.
\]

\important{Скорость сближения:}
\[
v_{\text{сбл}}=v_1+v_2.
\]

\subsection{Пример 1. Через сколько часов встретятся?}

Из двух городов, расстояние между которыми равно $560\ \km$, навстречу друг другу одновременно выехали два автомобиля со скоростями $60\ \kmh$ и $80\ \kmh$.

\begin{center}
\renewcommand{\arraystretch}{1.35}
\begin{tabularx}{0.7\linewidth}{X c c c}
\toprule
 & $v$ & $t$ & $S$ \\
\midrule
из $A$ & $60$ & $x$ & $60x$ \\
из $B$ & $80$ & $x$ & $80x$ \\
\bottomrule
\end{tabularx}
\end{center}

Составляем уравнение:
\[
60x+80x=560.
\]
\[
140x=560, \qquad x=4.
\]

\[
\boxed{4\ \ch}
\]

\subsection{Пример 2. На каком расстоянии от города A встретятся?}

Из городов $A$ и $B$, расстояние между которыми равно $480\ \km$, одновременно навстречу друг другу выехали два автомобиля. Из города $A$ автомобиль ехал со скоростью $55\ \kmh$, из города $B$ --- со скоростью $65\ \kmh$.

\begin{center}
\renewcommand{\arraystretch}{1.35}
\begin{tabularx}{0.7\linewidth}{X c c c}
\toprule
 & $v$ & $t$ & $S$ \\
\midrule
из $A$ & $55$ & $x$ & $55x$ \\
из $B$ & $65$ & $x$ & $65x$ \\
\bottomrule
\end{tabularx}
\end{center}

Сначала находим время до встречи:
\[
55x+65x=480.
\]
\[
120x=480, \qquad x=4.
\]

Расстояние от города $A$:
\[
S_A=55\cdot 4=220.
\]

\[
\boxed{220\ \km}
\]

\subsection{Тот же пример другим способом}

Иногда удобно обозначить за $x$ не время, а расстояние от города $A$ до места встречи.

Тогда:
\[
S_A=x, \qquad S_B=480-x.
\]

\begin{center}
\renewcommand{\arraystretch}{1.35}
\begin{tabularx}{0.8\linewidth}{X c c c}
\toprule
 & $v$ & $t$ & $S$ \\
\midrule
из $A$ & $55$ & $\dfrac{x}{55}$ & $x$ \\
из $B$ & $65$ & $\dfrac{480-x}{65}$ & $480-x$ \\
\bottomrule
\end{tabularx}
\end{center}

Так как автомобили выехали одновременно, времена равны:
\[
\frac{x}{55}=\frac{480-x}{65}.
\]

Решаем:
\[
65x=55(480-x),
\]
\[
65x=26400-55x,
\]
\[
120x=26400,
\]
\[
x=220.
\]

\[
\boxed{220\ \km}
\]

\subsection{Пример 3. Найти скорость}

Расстояние между городами равно $390\ \km$. Два автомобиля выехали одновременно навстречу друг другу. Скорость второго равна $60\ \kmh$, автомобили встретились через $3$ часа. Найдите скорость первого автомобиля.

\begin{center}
\renewcommand{\arraystretch}{1.35}
\begin{tabularx}{0.7\linewidth}{X c c c}
\toprule
 & $v$ & $t$ & $S$ \\
\midrule
первый & $x$ & $3$ & $3x$ \\
второй & $60$ & $3$ & $180$ \\
\bottomrule
\end{tabularx}
\end{center}

\[
3x+180=390.
\]
\[
3x=210, \qquad x=70.
\]

\[
\boxed{70\ \kmh}
\]

\section{Движение навстречу с задержкой старта}

\subsection{Ситуация}

Один объект выехал раньше, второй --- позже. Если второй ехал до встречи $x$ часов, а первый выехал на $h$ часов раньше, то первый ехал:
\[
x+h.
\]

Для движения навстречу уравнение обычно имеет вид:
\[
v_1(x+h)+v_2x=S.
\]

\important{Вопрос задачи важен:} если спрашивают время после выезда второго, ответом будет $x$, а не $x+h$.

\subsection{Пример 4. Найти время после выезда второго автомобиля}

Расстояние между городами $A$ и $B$ равно $580\ \km$. Из города $A$ в город $B$ со скоростью $80\ \kmh$ выехал автомобиль, а через $2$ часа после этого навстречу ему из города $B$ выехал второй автомобиль со скоростью $60\ \kmh$. Через сколько часов после выезда второго автомобиля они встретятся?

\begin{center}
\renewcommand{\arraystretch}{1.35}
\begin{tabularx}{0.75\linewidth}{X c c c}
\toprule
 & $v$ & $t$ & $S$ \\
\midrule
из $A$ & $80$ & $x+2$ & $80(x+2)$ \\
из $B$ & $60$ & $x$ & $60x$ \\
\bottomrule
\end{tabularx}
\end{center}

\[
80(x+2)+60x=580.
\]
\[
80x+160+60x=580.
\]
\[
140x=420.
\]
\[
x=3.
\]

\[
\boxed{3\ \ch}
\]

\subsection{Пример 5. Найти расстояние от города A}

Расстояние между городами $A$ и $B$ равно $380\ \km$. Из города $A$ в город $B$ со скоростью $50\ \kmh$ выехал автомобиль, а через $1$ час после этого навстречу ему из города $B$ выехал второй автомобиль со скоростью $60\ \kmh$. На каком расстоянии от города $A$ автомобили встретятся?

\begin{center}
\renewcommand{\arraystretch}{1.35}
\begin{tabularx}{0.75\linewidth}{X c c c}
\toprule
 & $v$ & $t$ & $S$ \\
\midrule
из $A$ & $50$ & $x+1$ & $50(x+1)$ \\
из $B$ & $60$ & $x$ & $60x$ \\
\bottomrule
\end{tabularx}
\end{center}

Сначала находим $x$ --- время движения второго автомобиля:
\[
50(x+1)+60x=380.
\]
\[
50x+50+60x=380.
\]
\[
110x=330.
\]
\[
x=3.
\]

Автомобиль из города $A$ ехал:
\[
x+1=4\ \ch.
\]

Расстояние от города $A$:
\[
S_A=50\cdot 4=200.
\]

\[
\boxed{200\ \km}
\]

\subsection{Пример 6. Найти скорость второго автомобиля}

Расстояние между городами $A$ и $B$ равно $440\ \km$. Из города $A$ в город $B$ со скоростью $60\ \kmh$ выехал автомобиль, а через $3$ часа после этого навстречу ему из города $B$ выехал второй автомобиль. Автомобили встретились через $2$ часа после выезда второго. Найдите скорость второго автомобиля.

Первый автомобиль ехал:
\[
3+2=5\ \ch.
\]

Он прошёл:
\[
60\cdot 5=300\ \km.
\]

Значит, второй автомобиль прошёл:
\[
440-300=140\ \km.
\]

Так как второй ехал $2$ часа, его скорость:
\[
v=\frac{140}{2}=70.
\]

\[
\boxed{70\ \kmh}
\]

\section{Движение вдогонку}

\subsection{Ситуация}

Два объекта движутся в одном направлении, но быстрый находится позади и догоняет медленного.

\catchupline{$AB$}{$v_{\text{быстр}}$}{$v_{\text{медл}}$}{движение вдогонку}

Если расстояние между ними сначала равно $D$, то быстрый должен не просто пройти свой путь, а ещё сократить отставание $D$.

Уравнение:
\[
v_{\text{быстр}}t=v_{\text{медл}}t+D.
\]

Или через скорость сближения:
\[
(v_{\text{быстр}}-v_{\text{медл}})t=D.
\]

\important{Скорость сближения при движении вдогонку:}
\[
v_{\text{сбл}}=v_{\text{быстр}}-v_{\text{медл}}.
\]

\subsection{Пример 7. Разность скоростей дана напрямую}

Города $A$, $B$ и $C$ соединены прямолинейным шоссе, город $B$ расположен между городами $A$ и $C$. Из города $A$ в сторону города $C$ выехал легковой автомобиль, и одновременно с ним из города $B$ в сторону города $C$ выехал грузовик. Скорость легкового автомобиля на $25\ \kmh$ больше скорости грузовика, расстояние $AB$ равно $125\ \km$.

\begin{center}
\renewcommand{\arraystretch}{1.35}
\begin{tabularx}{0.75\linewidth}{X c c c}
\toprule
 & $v$ & $t$ & $S$ \\
\midrule
легковой & $v+25$ & $x$ & $x(v+25)$ \\
грузовик & $v$ & $x$ & $xv$ \\
\bottomrule
\end{tabularx}
\end{center}

Легковой автомобиль должен пройти на $125\ \km$ больше:
\[
x(v+25)=xv+125.
\]
\[
xv+25x=xv+125.
\]
\[
25x=125.
\]
\[
x=5.
\]

\[
\boxed{5\ \ch}
\]

\subsection{Пример 8. Аналогичный случай}

Если разность скоростей равна $22\ \kmh$, а расстояние между городами $A$ и $B$ равно $132\ \km$, то:
\[
22x=132,
\]
\[
x=6.
\]

\[
\boxed{6\ \ch}
\]

\section{Расстояние между движущимися объектами}

\subsection{Ситуация}

Если два объекта стартуют из одного места и движутся в одном направлении, то расстояние между ними увеличивается за счёт разности скоростей:
\[
S_{\text{между}}=(v_1-v_2)t.
\]

\subsection{Пример 9. Пешеходы идут в одном направлении}

Два пешехода отправляются из одного места в одном направлении. Скорость первого на $2\ \kmh$ больше скорости второго. Через сколько минут расстояние между ними станет равным $300\ \m$?

Сначала переводим:
\[
300\ \m=0{,}3\ \km.
\]

\begin{center}
\renewcommand{\arraystretch}{1.35}
\begin{tabularx}{0.75\linewidth}{X c c c}
\toprule
 & $v$ & $t$ & $S$ \\
\midrule
быстрый & $y+2$ & $x$ & $x(y+2)$ \\
медленный & $y$ & $x$ & $xy$ \\
\bottomrule
\end{tabularx}
\end{center}

Расстояние между ними равно $0{,}3\ \km$:
\[
x(y+2)=xy+0{,}3.
\]
\[
xy+2x=xy+0{,}3.
\]
\[
2x=0{,}3.
\]
\[
x=0{,}15\ \ch.
\]

Переводим в минуты:
\[
0{,}15\cdot 60=9.
\]

\[
\boxed{9\ \text{минут}}
\]

\section{Движение до точки и обратно}

\subsection{Ситуация}

Два человека одновременно выходят из одного места к некоторой точке. Быстрый доходит до точки, разворачивается и идёт обратно. Нужно найти расстояние от точки разворота до места встречи.

\begin{center}
\begin{tikzpicture}[scale=0.95]
    \draw[line width=0.8pt] (0,0) -- (9,0);
    \fill[accent] (0,0) circle (2pt);
    \fill[accent] (9,0) circle (2pt);
    \fill[darkaccent] (6.8,0) circle (2pt);
    \node[below] at (0,0) {старт};
    \node[below] at (9,0) {опушка};
    \node[below] at (6.8,0) {встреча};
    \draw[->, line width=1pt, color=darkaccent] (0.6,0.45) -- (6.3,0.45);
    \draw[->, line width=1pt, color=darkaccent] (0.6,0.9) -- (8.6,0.9);
    \draw[->, line width=1pt, color=darkaccent] (8.6,0.9) -- (7.2,0.9);
    \node[above] at (3.3,0.45) {медленный};
    \node[above] at (5.2,0.9) {быстрый};
    \draw[<->, color=accent] (6.8,-0.55) -- (9,-0.55);
    \node[below] at (7.9,-0.55) {$x$};
\end{tikzpicture}
\end{center}

Если расстояние до опушки равно $L$, а $x$ --- расстояние от опушки до места встречи, то:
\[
S_{\text{медл}}=L-x,
\]
\[
S_{\text{быстр}}=L+x.
\]

Так как они вышли одновременно, времена равны:
\[
\frac{L-x}{v_{\text{медл}}}=\frac{L+x}{v_{\text{быстр}}}.
\]

\subsection{Пример 10. Найти расстояние от опушки}

Два человека отправляются из одного места на прогулку до опушки леса, находящейся в $6\ \km$ от места отправления. Один идёт со скоростью $3{,}5\ \kmh$, другой --- со скоростью $6{,}5\ \kmh$. Дойдя до опушки, второй с той же скоростью возвращается обратно. Сколько метров от опушки до места встречи?

Обозначим:
\[
x \text{ --- расстояние от опушки до места встречи в километрах}.
\]

\begin{center}
\renewcommand{\arraystretch}{1.35}
\begin{tabularx}{0.82\linewidth}{X c c c}
\toprule
 & $v$ & $t$ & $S$ \\
\midrule
медленный & $3{,}5$ & $\dfrac{6-x}{3{,}5}$ & $6-x$ \\
быстрый & $6{,}5$ & $\dfrac{6+x}{6{,}5}$ & $6+x$ \\
\bottomrule
\end{tabularx}
\end{center}

Времена равны:
\[
\frac{6-x}{3{,}5}=\frac{6+x}{6{,}5}.
\]

Решаем:
\[
6{,}5(6-x)=3{,}5(6+x).
\]
\[
39-6{,}5x=21+3{,}5x.
\]
\[
18=10x.
\]
\[
x=1{,}8\ \km.
\]

Переводим в метры:
\[
1{,}8\ \km=1800\ \m.
\]

\[
\boxed{1800\ \m}
\]

\subsection{Пример 11. Ещё один вариант}

Если скорости равны $4{,}5\ \kmh$ и $5{,}5\ \kmh$, а расстояние до опушки равно $6\ \km$, то:
\[
\frac{6-x}{4{,}5}=\frac{6+x}{5{,}5}.
\]
\[
5{,}5(6-x)=4{,}5(6+x).
\]
\[
33-5{,}5x=27+4{,}5x.
\]
\[
6=10x.
\]
\[
x=0{,}6\ \km=600\ \m.
\]

\[
\boxed{600\ \m}
\]

\section{Задачи с догонкой и разным временем прибытия}

\subsection{Ситуация}

Иногда в задаче есть две части:

\begin{enumerate}
    \item момент, когда один объект догоняет другой;
    \item момент, когда они приезжают в пункт назначения.
\end{enumerate}

В таких задачах удобно составлять две таблицы: одну для встречи, вторую для всего пути.

\subsection{Пример 12. Грузовик и легковой автомобиль}

Из города $A$ в город $B$ выехал грузовик, а через $1$ час следом за ним выехал легковой автомобиль. Через $2$ часа после выезда легковой автомобиль догнал грузовик и приехал в пункт $B$ на $3$ часа раньше, чем грузовик. Сколько часов потратил на дорогу от $A$ до $B$ грузовик?

Пусть:
\[
x \text{ --- скорость легкового автомобиля}, \qquad
y \text{ --- скорость грузовика}.
\]

До встречи легковой автомобиль ехал $2$ часа, а грузовик ехал $3$ часа.

\begin{center}
\renewcommand{\arraystretch}{1.35}
\begin{tabularx}{0.75\linewidth}{X c c c}
\toprule
 & $v$ & $t$ & $S$ \\
\midrule
легковой & $x$ & $2$ & $2x$ \\
грузовик & $y$ & $3$ & $3y$ \\
\bottomrule
\end{tabularx}
\end{center}

На момент встречи они прошли одинаковое расстояние:
\[
2x=3y.
\]

Пусть $t$ --- всё время движения грузовика от $A$ до $B$. Тогда легковой автомобиль был в пути на $3$ часа меньше:
\[
t-3.
\]

\begin{center}
\renewcommand{\arraystretch}{1.35}
\begin{tabularx}{0.75\linewidth}{X c c c}
\toprule
 & $v$ & $t$ & $S$ \\
\midrule
легковой & $x$ & $t-3$ & $x(t-3)$ \\
грузовик & $y$ & $t$ & $yt$ \\
\bottomrule
\end{tabularx}
\end{center}

Они проехали одно и то же расстояние $AB$:
\[
x(t-3)=yt.
\]

Из равенства $2x=3y$ получаем:
\[
\frac{x}{y}=\frac{3}{2}.
\]

Тогда:
\[
x(t-3)=yt
\]
делим на $y$:
\[
\frac{x}{y}(t-3)=t.
\]
Подставляем $\dfrac{x}{y}=\dfrac{3}{2}$:
\[
\frac{3}{2}(t-3)=t.
\]
\[
3(t-3)=2t.
\]
\[
3t-9=2t.
\]
\[
t=9.
\]

\[
\boxed{9\ \ch}
\]

\important{Важное замечание:} если в вопросе спрашивают, сколько часов потратил на дорогу грузовик, ответом является именно $t=9$, а не $9+3$. Величина $t$ уже обозначает всё время движения грузовика от $A$ до $B$.

\section{Алгоритм решения задач на движение}

\begin{enumerate}
    \item Определите тип движения: навстречу, вдогонку, из одной точки, туда-обратно.
    \item Выпишите известные скорости, времена и расстояния.
    \item Приведите единицы измерения к одному виду.
    \item Составьте таблицу с колонками $v$, $t$, $S$.
    \item Используйте формулу $S=vt$.
    \item Найдите, какие расстояния складываются или какие времена равны.
    \item Составьте уравнение.
    \item Решите уравнение.
    \item Проверьте, что ответ соответствует вопросу задачи.
\end{enumerate}

\section{Типичные ошибки}

\begin{enumerate}
    \item \important{Перепутать время.} Если второй автомобиль выехал позже, то первый ехал дольше.
    \item \important{Ответить не на тот вопрос.} В задаче могут спрашивать не время, а расстояние от города $A$.
    \item \important{Сложить скорости при движении вдогонку.} При догонке скорости вычитаются.
    \item \important{Не перевести метры в километры.} Например, $300\ \m=0{,}3\ \km$.
    \item \important{Неверно выбрать переменную.} Можно обозначать за $x$ время или расстояние, но таблица должна соответствовать выбранному смыслу.
    \item \important{Прибавить лишнее время к итоговому ответу.} Если $t$ уже обозначает всё время движения, прибавлять к нему ничего не нужно.
\end{enumerate}

\newpage

\section{Тренировочный блок}

Решите задачи. Ответы приведены на отдельной странице после тренировочного блока.

\subsection*{Средний уровень}

\begin{enumerate}
    \item Из двух городов, расстояние между которыми равно $420\ \km$, одновременно навстречу друг другу выехали два автомобиля со скоростями $50\ \kmh$ и $70\ \kmh$. Через сколько часов они встретятся?

    \item Из городов $A$ и $B$, расстояние между которыми равно $360\ \km$, одновременно навстречу друг другу выехали два автомобиля. Из города $A$ автомобиль ехал со скоростью $40\ \kmh$, из города $B$ --- со скоростью $50\ \kmh$. На каком расстоянии от города $A$ они встретятся?

    \item Расстояние между городами равно $450\ \km$. Два автомобиля выехали одновременно навстречу друг другу. Скорость первого равна $80\ \kmh$, автомобили встретились через $3$ часа. Найдите скорость второго автомобиля.
\end{enumerate}

\subsection*{Выше среднего уровня}

\begin{enumerate}[resume]
    \item Расстояние между городами $A$ и $B$ равно $520\ \km$. Из города $A$ в город $B$ выехал автомобиль со скоростью $70\ \kmh$, а через $2$ часа навстречу ему из города $B$ выехал второй автомобиль со скоростью $60\ \kmh$. Через сколько часов после выезда второго автомобиля они встретятся?

    \item Расстояние между городами $A$ и $B$ равно $390\ \km$. Из города $A$ в город $B$ выехал автомобиль со скоростью $45\ \kmh$, а через $1$ час навстречу ему из города $B$ выехал второй автомобиль со скоростью $60\ \kmh$. На каком расстоянии от города $A$ они встретятся?

    \item Города $A$, $B$ и $C$ расположены на одном шоссе, причём $B$ находится между $A$ и $C$. Из города $A$ в сторону $C$ выехал легковой автомобиль, одновременно из города $B$ в сторону $C$ выехал грузовик. Скорость легкового автомобиля на $18\ \kmh$ больше скорости грузовика, а расстояние $AB$ равно $126\ \km$. Через сколько часов легковой автомобиль догонит грузовик?

    \item Два пешехода вышли из одного места в одном направлении. Скорость первого на $3\ \kmh$ больше скорости второго. Через сколько минут расстояние между ними станет равным $450\ \m$?

    \item Два человека отправляются из одного места на прогулку до опушки леса, находящейся в $5\ \km$ от места отправления. Первый идёт со скоростью $4\ \kmh$, второй --- со скоростью $6\ \kmh$. Дойдя до опушки, второй с той же скоростью возвращается обратно. Сколько метров от опушки до места встречи?

    \item Расстояние между городами $A$ и $B$ равно $500\ \km$. Из города $A$ в город $B$ со скоростью $65\ \kmh$ выехал автомобиль, а через $2$ часа после этого навстречу ему из города $B$ выехал второй автомобиль. Они встретились через $4$ часа после выезда второго автомобиля. Найдите скорость второго автомобиля.
\end{enumerate}

\subsection*{Сложный уровень}

\begin{enumerate}[resume]
    \item Из города $A$ в город $B$ выехал грузовик, а через $1$ час следом за ним выехал легковой автомобиль. Через $3$ часа после своего выезда легковой автомобиль догнал грузовик и приехал в город $B$ на $2$ часа раньше грузовика. Сколько часов потратил на дорогу от $A$ до $B$ грузовик?
\end{enumerate}

\newpage

\section{Ответы к тренировочному блоку}

\begin{center}
\renewcommand{\arraystretch}{1.35}
\begin{tabularx}{0.9\linewidth}{c X}
\toprule
\textbf{№} & \textbf{Ответ} \\
\midrule
1 & $3{,}5\ \ch$ \\
2 & $160\ \km$ \\
3 & $70\ \kmh$ \\
4 & $2\ \ch$ \\
5 & $180\ \km$ \\
6 & $7\ \ch$ \\
7 & $9\ \text{минут}$ \\
8 & $1000\ \m$ \\
9 & $27{,}5\ \kmh$ \\
10 & $8\ \ch$ \\
\bottomrule
\end{tabularx}
\end{center}

\end{document}